\section{Intended Users of the Product}

SmartExam has four main user roles. A fifth role, SuperAdmin, manages institution-wide settings and administrator accounts.

\subsection{Administrator}
Administrators create, edit and deactivate user accounts, manage labs and workstations, and review device bindings and audit logs. They do not create exams or questions; that belongs to teachers.

\subsection{Teacher}
Teachers own the exam content and its logistics:
\begin{itemize}[noitemsep]
\item \textbf{Exams:} title, duration, schedule, MCQ, coding and theory questions, test cases, marks, allowed applications and lockdown mode.
\item \textbf{Eligibility:} import attendance, mark students eligible and assign each one to a lab workstation.
\item \textbf{After the exam:} review suggested marks, check the plagiarism matrix, override grades and download reports.
\end{itemize}

\subsection{Invigilator}
Invigilators watch a running exam on the live dashboard. Each student tile shows focus, active processes and a status of Normal, Warning, Suspicious or Violation. The status is set by the server's violation classifier. An invigilator can open a student's timeline and screenshots and send a warning to the student's screen. The invigilator is often the course teacher.

\subsection{Student}
Students only see the exams assigned to them. A student logs in on the bound workstation and passes the readiness check (network, clock, screen capture, webcam and microphone, secure mode, disk space). The student then starts the exam in the locked client. Inside it, the student can only read questions, write answers and submit. Practice exams let students rehearse beforehand without being graded. Once an exam is submitted, its answers can no longer be edited.
